% Nouveau EGA V (cote 25): archive cote 25, site page 45 (PDF page 46).
% typescript. Draft compared: none.
% AI-transcribed by Claude Opus 5.5 (claude-opus-5-5, adaptive thinking, effort high), Message Batches API, batch ega25b (msgbatch_01A7n5YuXubUPNK1zsepi8RG), 2026-11-13; not yet checked by eye.
\scanpage{25/45}
\mspage{\hand{\unsure{49|41}}}

hg$_1$ :  $g^{-1}(U) \longrightarrow$ PxT  n'est autre que l'inclusion canonique, puisqu'il
en est ainsi de sa restriction à $g^{-1}(V) \rightarrow$ V.  \var{\xout{Soit alors \unsure{U'} h$^{-1}$(}}{}{typewriter strike-out}{ega25_25_045-v0001}
\var{\xout{Comme h prend}}{}{typewriter strike-out}{ega25_25_045-v0002} Il en résulte que h \var{\xout{prend ses valeurs dans}}{}{typewriter strike-out}{ega25_25_045-v0003} se factorise
à travers l'adhérence schématique \var{\xout{de}}{}{typewriter strike-out}{ega25_25_045-v0004} $\overline{Y^1}$ de Y$^1$ dans  PxT , donc que
l'image inverse de Y$^1$ (qui est un ouvert de ladite adhérence) par h
est un ouvert de U, soit U' . En vertu de  h g$_1$ = inclusion, on
voit alors aussitôt que  g$_1$ induit un isomorphisme  $g^{-1}(U) \simeq$ U', et
\var{\xout{xxxxxx}}{}{typewriter strike-out}{ega25_25_045-v0005} cela montre bien que  $g^{-1}(U) \longrightarrow$ U est une immersion ouverte.
Lorsque X est propre sur k i.e. f ermé dans P, alors  g est propre donc
surjectif, donc  g$_1$ : $g^{-1}(U) \longrightarrow$ U est surjectif, donc un isomorphisme.
D'ailleurs, si  W est un ouvert de T tel que  $g^{-1}(W) \longrightarrow$ W soit un isomor%
phisme, il s'ensuit que W est lisse puisque Y$^1$ est lisse, donc W$\subset$U.
Cela achève la démonstration de 8.\var{\over{\unsure{8}}{9}}{9}{writing over}{ega25_25_045-v0006}. \var{\inserted{\unsure{$\underline{Y}^{\vee}$} \unsure{$r_{\text{sing}} = \rho$} \illegible{1}}}{\unsure{$\underline{Y}^{\vee}$} \unsure{$r_{\text{sing}} = \rho$} \illegible{1}}{insertion}{ega25_25_045-v0007}

\var{\xout{Remarque}}{}{typewriter strike-out}{ega25_25_045-v0008}  L\var{\over{'}{es}}{es}{writing over}{ega25_25_045-v0009} assertion\hand{s} finale\hand{s} de 8.\var{\over{\unsure{7}}{8}}{8}{writing over}{ega25_25_045-v0010}.\hand{,} dim $\underline{Y}^{\text{sup sing}}$ \var{\over{=}{\unsure{$\leq$}}}{\unsure{$\leq$}}{writing over}{ega25_25_045-v0011} r-2\hand{,} et
dim T = r-1 , sont triviales;  la première résult\var{\over{ant}{\unsure{e}}}{\unsure{e}}{writing over}{ega25_25_045-v0012} d\var{\over{u}{\unsure{e}}}{\unsure{e}}{writing over}{ega25_25_045-v0013} \var{\deleted{$\underline{Y}^{\text{sup sing}} \neq \underline{Y}^{\text{sing}}$}}{}{deletion}{ega25_25_045-v0014}
\var{\inserted{\unsure{$Y^{\text{sup sing}}$}}}{\unsure{$Y^{\text{sup sing}}$}}{insertion}{ega25_25_045-v0015} \var{\inserted{\unsure{contenu dans} $Y^{\text{sing}}$, \unsure{de dim r}}}{\unsure{contenu dans} $Y^{\text{sing}}$, \unsure{de dim r}}{insertion}{ega25_25_045-v0016} \var{\deleted{donc rare dans $\underline{Y}^{\text{sing}}$ puisque celui}}{}{deletion}{ega25_25_045-v0017} est irréductible, et du fait que
\var{\inserted{\unsure{est défini par l'annulation d'une section} D \unsure{d'un Module inversible}}}{\unsure{est défini par l'annulation d'une section} D \unsure{d'un Module inversible}}{insertion}{ega25_25_045-v0018} \var{\deleted{dim}}{}{deletion}{ega25_25_045-v0019} $\underline{Y}^{\text{sing}}$\hand{\unsure{sup sing}} \var{\deleted{= r-1}}{}{deletion}{ega25_25_045-v0020} ; la deuxième du fait que, L étant fini sur K, on a
deg tr L/k = deg\var{\xout{x}}{}{typewriter strike-out}{ega25_25_045-v0021} tr K/k, i.e.  \var{\xout{dim $\underline{Y}^{\text{sing}}$ xxxxxx}}{}{typewriter strike-out}{ega25_25_045-v0022} dim T = dim $\underline{Y}^{\text{sing}}$ = r-1.

\uline{Remarque} 8.\var{\over{\unsure{9}}{10}}{10}{writing over}{ega25_25_045-v0023}.  Comme on a si\var{\over{n}{g}}{g}{writing over}{ega25_25_045-v0024}\unsure{n}alé dans  8.\var{\over{\unsure{8}}{9}}{9}{writing over}{ega25_25_045-v0025}. , avec les notations de
ce corollaire, on a  $g^{-1}(U) \subset \underline{Y}^{\text{sing}} - \underline{Y}^{\text{sup sing}}$ ; mais on notera que
même si X est fermé dans P,  cette inclusion n'est pas nécessairement
\var{\margin{\unsure{$\ast$}}}{\placedmargin{\unsure{$\ast$}}}{marginal note placed}{ega25_25_045-v0026} une égalité, en d'autres termes (notant que $g^{-1}(U)$ n'est autre que
l'ensemble des points où g est \uline{étale}, alors que $\underline{Y}^{\text{sing}}$-$\underline{Y}^{\text{sup sing}}$ est
lèensemble des points où g est \uline{non ramifié}) il peut y avoir des points
y de $\underline{Y}^{\text{sing}}$ en \hand{\unsure{´}}lesquels  g est \var{\xout{étale}}{}{typewriter strike-out}{ega25_25_045-v0027} non ramifié, mais non étale (ce qui
\var{\inserted{\var{\deleted{\illegible{1}}}{}{deletion}{ega25_25_045-v0029}}}{\var{\deleted{\illegible{1}}}{}{deletion}{ega25_25_045-v0029}}{insertion}{ega25_25_045-v0028} implique d'ailleurs que g(y) est un point non \var{\xout{normal}}{}{typewriter strike-out}{ega25_25_045-v0030} géom normal et même
non géom unibranche de  T)\hand{.} \var{\xout{xxxxxx}}{}{typewriter strike-out}{ega25_25_045-v0031} En termes géométriques, cela corres%
pond au phénomène suivant: on peut avoir un hyperplan tangent non oscul
% NOTE: the last line is cut at the right edge of the scan
\var{\margin{45}}{\placedmargin{45}}{marginal note placed}{ega25_25_045-v0032}
